\documentclass{article}
\usepackage[margin=1in]{geometry}
\usepackage{nockmark}
\usepackage{listings}
\lstset{basicstyle=\ttfamily\small,columns=fullflexible}
\title{\texttt{nockmark}: annotated code fragments}
\author{Sigilante}
\date{v0.1.0, 2026-11-06}
\begin{document}
\maketitle

\texttt{nockmark} draws a code fragment as a column of lines with boxed spans
and notes, using TikZ.  It was made for Nock, which has no layout of its own,
but any monospaced code works.  It runs under pdf\LaTeX, Xe\LaTeX\ and
Lua\LaTeX.

\section{Example}

\begin{center}
\begin{nockmark}
  \nkheading{Static sites: the code travels with the literal}
  \nkline{\nk{[8 [1 0]}}
  \nkline[1]{\nk{8 }\nkstatic{s1}{[1 \nkmeta{B}]}}
  \nkline[1]{\nkstatic{s2}{9 2 0 1}\nk{]}}
  \nknote{s1}{literal formula:  a \texttt{Form} carrying \texttt{nockToSki} \textit{B}}
  \nknote{s2}{opcode~9 pulls the \texttt{Form}:  \texttt{app} runs its code}
  \nkskip{0.3}
  \nkline{\nk{\nkmeta{B} = [6 [5 [0 7] 4 0 6] [0 6]}}
  \nkline[4]{\nkstatic{s3}{9 2 [0 2] [4 0 6] 0 7}\nk{]}}
  \nknotebelow{s3}{the arm is copied from the subject by \texttt{[0 2]}:  still the \texttt{Form}}
  \nkrule
  \nkheading{Dynamic site: a formula built from data}
  \nkline{\nk{[2 [0 1] }\nkdynamic{d1}{[1 4] [1 0 1]}\nk{]}}
  \nknotebelow{d1}{autocons builds the plain \texttt{Cell} \texttt{[4 0 1]}:  \texttt{app} hands it to \texttt{nockI}}
\end{nockmark}
\end{center}

\noindent The source:
\begin{lstlisting}
\begin{nockmark}
  \nkheading{Static sites: the code travels with the literal}
  \nkline{\nk{[8 [1 0]}}
  \nkline[1]{\nk{8 }\nkstatic{s1}{[1 \nkmeta{B}]}}
  \nkline[1]{\nkstatic{s2}{9 2 0 1}\nk{]}}
  \nknote{s1}{literal formula: a \texttt{Form} carrying ...}
  \nknote{s2}{opcode~9 pulls the \texttt{Form}: ...}
  \nkskip{0.3}
  \nkline{\nk{\nkmeta{B} = [6 [5 [0 7] 4 0 6] [0 6]}}
  \nkline[4]{\nkstatic{s3}{9 2 [0 2] [4 0 6] 0 7}\nk{]}}
  \nknotebelow{s3}{the arm is copied from the subject ...}
  \nkrule
  \nkheading{Dynamic site: a formula built from data}
  \nkline{\nk{[2 [0 1] }\nkdynamic{d1}{[1 4] [1 0 1]}\nk{]}}
  \nknotebelow{d1}{autocons builds the plain \texttt{Cell} ...}
\end{nockmark}
\end{lstlisting}

\section{Commands}

\begin{description}
\item[\texttt{\textbackslash begin\{nockmark\}[\emph{keys}]}] a fragment; it is a \texttt{tikzpicture}, so place it in a \texttt{figure} or \texttt{center}.
\item[\texttt{\textbackslash nkline[\emph{n}]\{\emph{pieces}\}}] a line at the next baseline, indented $n$ steps.
\item[\texttt{\textbackslash nk\{\emph{code}\}}] plain code.  A trailing space is kept.
\item[\texttt{\textbackslash nkstatic\{\emph{name}\}\{\emph{code}\}}] a solid box, named for notes.
\item[\texttt{\textbackslash nkdynamic\{\emph{name}\}\{\emph{code}\}}] a dashed box.
\item[\texttt{\textbackslash nkbox[\emph{style}]\{\emph{name}\}\{\emph{code}\}}] a box in any TikZ style, e.g.\ one defined with \texttt{\textbackslash tikzset}.
\item[\texttt{\textbackslash nkmeta\{\emph{x}\}}] a metavariable inside code, in the text face.
\item[\texttt{\textbackslash nknote[\emph{opts}]\{\emph{name}\}\{\emph{text}\}}] a note in the right column, with an arrow to the end of the box's line.
\item[\texttt{\textbackslash nknotebelow[\emph{opts}]\{\emph{name}\}\{\emph{text}\}}] a note under the box; following lines clear it.
\item[\texttt{\textbackslash nkheading\{\emph{text}\}}, \texttt{\textbackslash nkrule}, \texttt{\textbackslash nkskip\{\emph{cm}\}}] a heading, a dotted divider, extra space.
\item[\texttt{\textbackslash nockmarkset\{\emph{keys}\}}] set keys for every later fragment.
\end{description}

\section{Keys}

Lengths without a unit are in centimetres.  Defaults in parentheses.

\begin{description}
\item[\texttt{line sep}] (0.6) baseline to baseline.
\item[\texttt{indent}] (0.25) one indent step.
\item[\texttt{heading sep}] (0.75) heading baseline to the next line.
\item[\texttt{note column}] (2.9) left edge of right-column notes.
\item[\texttt{note width}] (4.9cm), \textbf{\texttt{below width}} (7.2cm) note widths.
\item[\texttt{below skip}] (0.45) from a below-note's bottom to the next baseline.
\item[\texttt{rule raise}] (0.15), \textbf{\texttt{rule skip}} (0.4) rule placement.
\item[\texttt{box gap}] (1.6pt) box padding and spacing.
\item[\texttt{code font}, \texttt{note font}, \texttt{heading font}, \texttt{meta font}] the faces.
\end{description}

\noindent The TikZ styles \texttt{nk code}, \texttt{nk box}, \texttt{nk static},
\texttt{nk dynamic}, \texttt{nk note}, \texttt{nk heading}, \texttt{nk arrow} and
\texttt{nk rule} can be redefined with \texttt{\textbackslash tikzset}.

\section{Other code}

\begin{center}
\begin{nockmark}[note column=5.2, note width=5cm]
  \nkline{\nk{++\ \ dec}}
  \nkline[1]{\nk{|=\ \ a=@}}
  \nkline[1]{\nkstatic{h}{?<\ \ =(0 a)\ \ !!}}
  \nknote{h}{crash on zero}
  \nkline[1]{\nk{=+\ \ b=0}}
  \nkline[1]{\nk{|-\ \ ?:\ \ }\nkdynamic{l}{=(a +(b))}\nk{\ \ b}}
  \nknote{l}{the loop test}
  \nkline[1]{\nk{\$(b +(b))}}
\end{nockmark}
\end{center}

\end{document}
